% PAGES: 204-207
% Continues after sec06_c4_1.tex (PDF pages 201-203), which ended section
% 6.5 "References" cleanly with no environment left open. PDF page 204
% (folio 188) opens fresh with the "6.6 Exercises" section heading at the
% top of the page.
%
% This is the whole of Chapter 6's exercise list (items 1-13, PDF pages
% 204-207 / folios 188-191), kept in a single \exercises environment so the
% shared enumerate counter is never reset mid-list. PDF page 207 (folio
% 191) ends with the note to exercise 13; the remainder of that page is
% blank. PDF page 208 (folio 192) was rendered and inspected directly and
% contains no printed text at all (genuinely blank), confirming the end of
% Chapter 6. No environment is left open at the end of this file, and
% nothing further follows in this chapter.

\section{Exercises}

\begin{exercises}

\item Let
\[
A \;=\; \begin{pmatrix} 2 & -2 \\ -2 & 5 \end{pmatrix}.
\]
\begin{subparts}
\item Find a $2 \times 2$ matrix $M$ such that $A = M^2$;
\item Find a $2 \times 2$ matrix $M$ such that $A = MM\tp$.
\end{subparts}

\item Let $A = \begin{pmatrix} 41 & 12 \\ 12 & 34 \end{pmatrix}$.
\begin{subparts}
\item Show that the matrix $A$ is symmetric positive definite.
\item Find a $2 \times 2$ matrix $B$ such that $A = B^2$.
\item Find a $2 \times 2$ upper triangular matrix $U$ such that $A = U\tp U$.
\end{subparts}

\item Let $A = \begin{pmatrix} 1 & -0.2 & 0.8 \\ -0.2 & 2 & 0.3 \\ 0.8 & 0.3 & 1.1
\end{pmatrix}$.
\begin{subparts}
\item Show that the matrix $A$ is weakly diagonally dominated, and therefore symmetric
positive semidefinite.
\item Use Sylvester's criterion to show that the matrix $A$ is symmetric positive
definite.
\end{subparts}

\item Let $A$ be a symmetric positive definite matrix, and let $U$ be the Cholesky
factor of $A$. If $V$ is an upper triangular matrix such that $A = V\tp V$, show that
there exists a diagonal matrix $D$ whose entries on the main diagonal are either $-1$
or $1$ such that $V = DU$.

\item Let $B_N$ be the following $N \times N$ symmetric positive definite matrix:
\begin{equation}
B_N \;=\;
\begin{pmatrix}
2 & -1 & \dots & 0 \\
-1 & \ddots & \ddots & \vdots \\
\vdots & \ddots & \ddots & -1 \\
0 & \dots & -1 & 2
\end{pmatrix}.
\end{equation}

Let $C_N$ be an $N \times (N+1)$ matrix with
\begin{align*}
C_N(i,i) &= 1, \quad \forall\, i = 1:N; \\
C_N(i,i+1) &= -1, \quad \forall\, i = 1:N,
\end{align*}
and with all the other entries equal to 0.

Show that $B_N = C_N C_N\tp$. Why is this not the Cholesky decomposition of the matrix
$B_N$?

\item Let $B_N$ be the $N \times N$ tridiagonal symmetric positive definite matrix
given by (6.88). Show that the Cholesky factor $U_N$ of the matrix $B_N$ is the upper
triangular bidiagonal matrix given by
\begin{equation}
U_N(i,i) \;=\; \sqrt{\frac{i+1}{i}}, \quad \forall\, i = 1:N;
\end{equation}
\begin{equation}
U_N(i,i+1) \;=\; -\sqrt{\frac{i}{i+1}}, \quad \forall\, i = 1:(N-1).
\end{equation}

\item Let $B_N$ be the $N \times N$ tridiagonal symmetric positive definite matrix
given by (6.88), and let $U_N$ be the Cholesky factor of $B_N$ given by (6.89--6.90).

\begin{subparts}
\item Show that the solution to a linear system $B_N x = b$, where $b$ and $x$ are
$N \times 1$ column vectors can be obtained by using the explicit pseudocode below:

\begin{center}
\fbox{%
\begin{minipage}{0.85\linewidth}
Function Call: \\
$x = \text{linear\_solve\_cholesky\_B\_N}(b)$
\\[0.8em]
Input: \\
$b = N \times 1$ column vector
\\[0.8em]
Output: \\
$x =$ solution to $B_N x = b$
\\[0.8em]
$y(1) = \dfrac{b(1)}{\sqrt{2}}$ \\
for $i = 2:N$ \\
\hspace*{1.5em}$y(i) = \dfrac{b(i) + y(i-1)\sqrt{(i-1)/i}}{\sqrt{(i+1)/i}}$ \\
end \\
$x(N) = \dfrac{y(N)\sqrt{N}}{\sqrt{N+1}}$ \\
for $i = (N-1):1$ \\
\hspace*{1.5em}$x(i) = \dfrac{y(i) + x(i+1)\sqrt{i/(i+1)}}{\sqrt{(i+1)/i}}$ \\
end
\end{minipage}%
}
\end{center}

\item What is the operation count for the pseudocode above, and how does it compare to
$8n + O(1)$, the operation count for the optimal linear solver for tridiagonal
symmetric positive definite matrices?
\end{subparts}

\item Let $B_N$ be the $N \times N$ tridiagonal symmetric positive definite matrix
given by (6.88). Show that the LU factors $L$ and $U$ of the matrix $B_N$ are the
lower triangular bidiagonal matrix and the upper triangular bidiagonal matrix given by
\begin{equation}
L(i,i) = 1, \quad \forall\, i = 1:N; \quad L(i+1,i) = -\frac{i}{i+1}, \quad \forall\, i
= 1:(N-1);
\end{equation}
\begin{equation}
U(i,i) = \frac{i+1}{i}, \quad \forall\, i = 1:N; \quad U(i,i+1) = -1, \quad \forall\, i
= 1:(N-1).
\end{equation}

\item Let $B_N$ be the $N \times N$ tridiagonal symmetric positive definite matrix
given by (6.88), and let $L$ and $U$ be the LU factors of $B_N$ given by (6.91--6.92).

\begin{subparts}
\item Show that the solution to a linear system $B_N x = b$, where $b$ and $x$ are
$N \times 1$ column vectors can be obtained by using the explicit pseudocode below:

\begin{center}
\fbox{%
\begin{minipage}{0.85\linewidth}
Function Call: \\
$x = \text{linear\_solve\_lu\_B\_N}(b)$
\\[0.8em]
Input: \\
$b = N \times 1$ column vector
\\[0.8em]
Output: \\
$x =$ solution to $B_N x = b$
\\[0.8em]
$y(1) = b(1)$ \\
for $i = 2:N$ \\
\hspace*{1.5em}$y(i) = b(i) + \dfrac{(i-1)y(i-1)}{i}$ \\
end \\
$x(N) = \dfrac{Ny(N)}{N+1}$ \\
for $i = (N-1):1$ \\
\hspace*{1.5em}$x(i) = \dfrac{i(y(i)+x(i+1))}{i+1}$ \\
end
\end{minipage}%
}
\end{center}

\item What is the operation count for the pseudocode above, and how does it compare to
$8n + O(1)$, the operation count for the optimal linear solver for tridiagonal
symmetric positive definite matrices?
\end{subparts}

\item Write an explicit optimal pseudocode for solving linear systems corresponding to
the same tridiagonal symmetric positive definite matrix. In other words, write a
pseudocode for solving $p$ linear systems $Ax_i = b_i$, for $i = 1:p$, where $A$ is an
$n \times n$ tridiagonal symmetric positive definite matrix, by using the explicit
solver linear\_solve\_LU\_tridiag\_spd from Table 6.6 and the method from Table 6.3.
What is the operation count for solving the $p$ linear systems?

\item Write the pseudocode for the Cholesky decomposition of symmetric positive
definite banded matrices of band $m$. What is the corresponding operation count?

\noindent Hint: Use (without proving) the fact that the Cholesky factor of a symmetric
positive definite banded matrix of band $m$ is a banded upper triangular matrix of band
$m$.

\item What is the operation count for solving a linear system corresponding to a
symmetric positive definite banded matrix of band $m$ using a Cholesky linear solver?

\item The following discount factors were obtained from market data:
\begin{center}
\begin{tabular}{cc}
Date & Discount Factor \\[0.4em]
2 months & 0.9980 \\
5 months & 0.9935 \\
11 months & 0.9820 \\
15 months & 0.9775
\end{tabular}
\end{center}

The overnight rate is $0.75\%$.

\begin{subparts}
\item What are the corresponding 2 months, 5 months, 11 months and 15 months zero
rates?
\item What is the tridiagonal system that must be solved in the efficient
implementation of the natural cubic spline interpolation for finding the zero rate
curve for all times less than 15 months?
\item Use the efficient implementation of the natural cubic spline interpolation to
find a zero rate curve for all times less than 15 months matching the discount factors
above.
\item Find the value of a 14 months quarterly coupon bond with $2.5\%$ coupon rate.
\end{subparts}

\noindent Note: A quarterly coupon bond with face value \$100, coupon rate $C$, and
maturity $T$ pays the holder of the bond a coupon payment equal to $\frac{C}{4} \cdot
100$ every three months, except at maturity. The final payment at maturity $T$ is equal
to the face value of the bond plus one coupon payment, i.e., $100 + \frac{C}{4}100$.

\end{exercises}
